\subsection*{pH probe calibration methodology}
% Text suitable for the Capstone Methodology chapter (SEN0169 software two-point calibration).
% Session documented: 16 September 2026, approximately 19:19--19:27 AEST.
% Supporting data: \texttt{ph\_calibration\_timeseries.csv}, \texttt{ph\_calibration\_events.csv}.

\subsubsection{Hardware and signal path}
A DFRobot SEN0169 industrial pH probe and meter board were used to measure
reservoir acidity. The meter board outputs an analogue voltage in the
$0$--$5\,\mathrm{V}$ range. This signal was scaled into the ESP32 ADC input
range with a discrete resistive divider ($1\,\mathrm{k}\Omega$ /
$2\,\mathrm{k}\Omega$), so the firmware reconstructs board voltage as
\begin{equation}
  V_{\mathrm{board}} = V_{\mathrm{GPIO}} \times 1.5.
\end{equation}
Converted $\mathrm{pH}$ is obtained from a linear model
\begin{equation}
  \mathrm{pH} = m\,V_{\mathrm{board}} + b,
\end{equation}
where the slope $m$ and intercept $b$ are determined by software two-point
calibration against standard buffers and stored in on-device flash
(LittleFS). Until a successful calibration is saved, the node uses the
DFRobot library factory defaults
($m=-5.6344$, $b=15.4516$).

\subsubsection{Calibration procedure}
Calibration was performed on 16~September~2026 using the Streamlit dashboard
wizard, which publishes MQTT commands
(\texttt{\{"cmd":"ph\_cal","action":\ldots\}}) to the ESP32. The operator
procedure was:
\begin{enumerate}
  \item Rinse the probe in distilled water and blot dry.
  \item Place the probe in $\mathrm{pH}\,4.00$ buffer and wait for
        $V_{\mathrm{board}}$ to settle.
  \item Capture the acid point via the dashboard (\emph{Capture acid}).
  \item Rinse again, then place the probe in $\mathrm{pH}\,9.18$ buffer and
        wait for settle.
  \item Capture the alkali point (\emph{Capture alkali}).
  \item Save the fit (\emph{Save \& exit}), which writes $m$ and $b$ to
        LittleFS and resumes normal telemetry.
  \item Return the probe to distilled rinse water as a qualitative check
        (distilled water is expected to read slightly acidic).
\end{enumerate}

From the logged session, acid capture occurred at
\textbf{19:19:46\,AEST} with
$V_a = 1.820\,\mathrm{V}$ assigned to $\mathrm{pH}\,4.00$, and alkali
capture at \textbf{19:23:34\,AEST} with
$V_b = 2.157\,\mathrm{V}$ assigned to $\mathrm{pH}\,9.18$.
The firmware fitted the unique line through these two points
(no imposed slope polarity):
\begin{equation}
  m = \frac{9.18-4.00}{V_b-V_a},\qquad
  b = 4.00 - m\,V_a,
\end{equation}
yielding $m \approx 15.370$ and $b \approx -23.970$, saved at
\textbf{19:24:40\,AEST}. After save, with the probe in distilled rinse water
(\(\approx\)19:27\,AEST), reported $\mathrm{pH}$ was in the mid-$5$ range,
consistent with slightly acidic rinse water.

\subsubsection{Data logging used for documentation}
Continuous sensor samples were retrieved from InfluxDB Cloud
(\texttt{sensor\_reading} measurement, device \texttt{garden-node-01}) for
the calibration window. Calibration state markers
(\texttt{ph\_cal\_message}, pending voltages, fitted $m$ and $b$) were
aligned from concurrent MQTT telemetry archived by the local dashboard.
The combined export
\texttt{ph\_calibration\_timeseries.csv} and event summary
\texttt{ph\_calibration\_events.csv} were used to generate time-series
figures of $V_{\mathrm{board}}$ and reported $\mathrm{pH}$, with vertical
markers at acid capture, alkali capture, save, and the distilled-water
check.

\subsubsection{Notes for interpretation}
The two captured voltages in this session produced a positive slope
($V_b>V_a$ while $\mathrm{pH}_b>\mathrm{pH}_a$). The software intentionally
accepts the empirical two-point line so that the documented operator
captures are exactly reproducible. Independent strip/reference checks of
tank samples remain useful to validate that buffer order and probe settle
time were correct before relying on the fitted line for closed-loop dosing.
